% !TeX root = ../偏微分方程笔记.tex
% 由主文档通过 \input 载入；请编译主文档。
\setcounter{chapter}{0} % 保留原书章号。
\chapter{偏微分方程的基本概念}
\label{chap:concepts}

偏微分方程把一个场在不同方向的变化联系起来。例如温度是位置与时间的函数，温度的时间变化又受到空间温差的影响。因此，首先应辨认未知量、独立变量和定义域，再讨论方程的结构。

\section{从常微分方程到偏微分方程}
\begin{definition}
设 $\Omega\subset\mathbb R^n$ 为开集。含有未知函数 $u:\Omega\to\mathbb R$ 及其偏导数的关系称为偏微分方程；其中实际出现的最高阶偏导数的阶数称为方程的\emph{阶（order）}。二阶标量方程可写为
\[
 F(x,u,\nabla u,D^2u)=0,
\]
其中 $D^2u=(\partial_{x_i x_j}u)_{i,j=1}^n$ 为海森矩阵（Hessian matrix），$F$ 为已知函数。
\end{definition}
当未知量是向量时得到方程组。时间也是独立变量，但为突出演化方向通常写成 $u(x,t)$。热方程 $u_t-\kappa u_{xx}=0$ 按所有导数的最高阶计是二阶方程，按时间导数看则是一阶演化方程。

在矩形 $I\times J$ 上，$u_x=0$ 的所有 $C^1$ 解为 $u(x,y)=g(y)$，其中 $g$ 可任取 $C^1$ 函数：固定 $y$ 后，一元函数 $x\mapsto u(x,y)$ 的导数为零，故沿水平线为常数。偏微分方程的解通常包含函数自由度，仅写方程不足以确定解。定义在某点邻域内的解称局部解；定义在所讨论的整个区域或全部正时间上的解称整体解，二者不能混同。

这里的“任意函数”与常微分方程中的“任意常数”有本质区别。求 $y'(x)=0$，只需给出一个数 $y(0)$；求 $u_x=0$，即使已知 $u(0,0)$，仍不知道其他水平线上的常数是多少，必须给出整条横截线上的数据 $u(0,y)$。这也是偏微分方程的附加条件通常是函数的原因。另一方面，偏导数是在其余变量固定时求导：若 $u(x,t)=e^{-t}\sin x$，则 $u_t=-e^{-t}\sin x$、$u_{xx}=-e^{-t}\sin x$，所以它满足 $u_t=u_{xx}$；沿运动点 $x=x(t)$ 观测时，则有
\[
 \frac d{dt}u(x(t),t)=u_t(x(t),t)+x'(t)u_x(x(t),t).
\]
这个全导数一般不等于 $u_t$，却正是特征线法所要利用的量。

\section{线性结构与非线性的层次}
\begin{definition}
二阶方程若能写成
\[
 Lu:=\sum_{i,j=1}^n a_{ij}(x)u_{x_ix_j}
       +\sum_{i=1}^n b_i(x)u_{x_i}+d(x)u=f(x),
\]
则称为\emph{线性方程（linear equation）}，因为 $L(\alpha u+\beta v)=\alpha Lu+\beta Lv$。$f=0$ 时称齐次，否则称非齐次。若最高阶项仍线性且其系数只依赖 $x$，低阶项允许非线性，称\emph{半线性（semilinear）}；若最高阶项线性而系数还可依赖 $u,\nabla u$，称\emph{拟线性（quasilinear）}；最高阶导数以非线性方式出现时称\emph{完全非线性（fully nonlinear）}。
\end{definition}
这是一组逐层包含的结构类别，具体分类时通常指出最精细的一类。一阶方程的拟线性系数可依赖 $u$；二阶方程的拟线性系数还可依赖一阶导数。

\begin{example}
$u_t-u_{xx}=u^3$ 是二阶半线性方程；$u_t+u u_x=0$ 是一阶拟线性方程；$u_{xx}u_{yy}-u_{xy}^2=1$ 是二阶完全非线性方程。$x^2u_{xx}+\sin x\,u_y=0$ 仍是线性方程，自变量进入系数不破坏线性。
\end{example}
若 $Lu_1=Lu_2=f$，则 $L(u_1-u_2)=0$。因而非齐次方程的解集在非空时是“一个特解加齐次解空间”。但叠加还必须与边界数据匹配：两个具有相同非零边值的解相加，一般不再具有原边值。

判断线性时，应把自变量与未知量分开看。$x u_{xx}$ 中的 $x$ 已知，乘法只改变系数；$u u_{xx}$ 中的第一个 $u$ 也是未知量，因此产生非线性。还应把方程和边界条件分别分类：$\Delta u=0$ 配上 $\partial_\nu u+u^3=0$，微分方程线性，但整个边值问题非线性。反过来，非零边界数据并不破坏线性算子结构，只是使允许的函数集合不再经过零函数。

\section{经典解、广义解与解的验证}
\begin{definition}
若未知函数具有方程中要求的连续偏导数，并在定义域内逐点满足方程，就称为\emph{经典解（classical solution）}。对热方程常写 $u\in C^{2,1}$，表示所需的二阶空间导数和一阶时间导数连续。完整定解问题还要求在指定意义下满足初始条件和边界条件。
\end{definition}
验证解有三个环节：检查定义域与可微性，代入微分方程，核对全部附加数据。例如
\[
 u(x,t)=e^{-\kappa(\pi/\ell)^2t}\sin(\pi x/\ell)
\]
满足 $u_t=\kappa u_{xx}$、$u(0,t)=u(\ell,t)=0$，其初值只能是 $\sin(\pi x/\ell)$，不能据此宣称它解出了任意初温问题。

若初值 $g(x)=|x|$，输运表达式 $u(x,t)=|x-at|$ 沿直线 $x=at$ 不可微，不能作为该处的经典解。将导数转移到光滑测试函数后，它仍可满足积分恒等式。这引出\emph{弱解（weak solution）}：保留方程的积分意义，放宽解本身的可微性。严格定义见第\ref{chap:sobolev}章。

\section{适定性与研究问题的层次}
\begin{definition}
设数据属于赋范空间 $Y$，解属于赋范空间 $X$。若每个允许的数据 $d\in Y$ 都有唯一解 $u\in X$，且数据到解的映射连续，则称问题在这些空间中\emph{适定（well-posed）}。连续依赖的含义是 $d_m\to d$ 蕴含相应的 $u_m\to u$，分别采用 $Y$ 和 $X$ 的范数。
\end{definition}
存在性回答能否求解，唯一性回答是否由数据确定，稳定性回答数据误差会怎样影响结果。显式公式是求解工具，适定性则是问题本身的性质；许多适定问题没有可用的闭式解。逆向热方程有些数据确实能产生解，但高频误差会剧烈放大，因此连续依赖仍可能失败，见第\ref{chap:heat}章。

\begin{example}[用最简单的问题区分存在性与唯一性]
在 $0<x<1$ 上考虑 $-u''=0$。若给定 $u(0)=A,u(1)=B$，通解 $u=ax+b$ 中的两个常数被确定，得到唯一解 $u=A+(B-A)x$。若只给 $u(0)=A$，斜率仍任意；若改给 $u'(0)=0,u'(1)=1$，由于 $u'$ 应为常数，问题无解；若给 $u'(0)=u'(1)=0$，则任意常数都是解。附加条件的数量相同，也不意味着可解性相同，条件之间的关系同样重要。
\end{example}

这个例子虽然只有一个空间变量，却已经展示了椭圆边值问题的主要问题。学习多维理论时，可以先把抽象结论放回区间模型检验。稳定性则需要说明怎样量误差：$L^2$ 范数衡量平均平方误差，上确界范数衡量最大的点值误差；在一种范数中稳定，不自动说明在另一种范数中也稳定。

\section{例题与习题}
\begin{example}
在矩形上求 $u_{xy}=0$ 的全部 $C^2$ 解，并附加 $u(x,0)=p(x)$、$u(0,y)=q(y)$。
\end{example}
\begin{solve}
固定 $x$，$u_x$ 与 $y$ 无关，所以积分得 $u(x,y)=A(x)+B(y)$。两条坐标轴相交处必须满足 $p(0)=q(0)$；在此前提下唯一解为
\[
 u(x,y)=p(x)+q(y)-p(0).
\]
其中 $p,q\in C^2$，且矩形含有相应轴段。代入即可验证；若交点数据不同，则不存在连续解。
\end{solve}
\begin{enumerate}
 \item 分类 $u_t=(1+u^2)u_{xx}$ 与 $u_t=u_{xx}+(u_x)^2$，说明判别依据。
 \item 证明线性齐次边值问题的解集是线性空间，非齐次问题的非空解集通常只是仿射空间。
 \item 对 $u_t+a u_x=0$ 的平移解，证明初值在上确界范数中的误差在演化中保持不变。
\end{enumerate}

